% Pursuing Stacks (cotes 134-2 to 134-8): archive cote 134-3, site page 59 (PDF page 60), typescript page 59.
% typescript. Draft compared: ps2.tex:2991-3069.
% AI-transcribed by Claude Opus 5.5 (claude-opus-5-5, adaptive thinking, effort high), Message Batches API, batch ps1 (msgbatch_01Rv5hCaNTtUt4JHx1yqoHPL), 2026-11-12; not yet checked by eye.
\scanpage{134-3/59}
\mspage{59}

On a more technical level,  I got aware too that the last corollary
\margin{\unsure{L} $(^{42})$} stated was slightly uncorrect,  because it is by no means clear that the
crible \inserted{$C_u$} constructed there is  functorial in u  i.e. corresponds to a map
I\hand{$\to\alpha$}(\hand{$\Delta$}$_1$) ,  this in fact has still to be assum\xout{m}ed  (and turns out to be
\xout{automatic} satisfied in all cases which turned up naturally so far and which
I looked up).  But let's now \hand{``}retell the story\hand{''}!

One key notion visibly in the homotopy technique used,  and which
needs a name in the long last,  is the notion of a \uline{homotopy} \uline{interval}\xout{\unsure{,}}
( "segment homotopique" in French).  To be really outspoken about the very
formal nature of this notion and the way it will be used,  let's develop
it in any category \hand{$\mathcal{A}$} endowed with a subset  W \hand{$\subset$} Fl(\hand{$\mathcal{A}$}) of the set of
arrows of \hand{$\mathcal{A}$} , and satisfying the usual conditions :

a)  W contains all isomorphisms of \hand{$\mathcal{A}$} ,

b)  for two composable arrows u,v in \hand{$\mathcal{A}$},  if  two among u,v, vu are
in W,  so is the third,  and

c)  if  i: $F_o$ \hand{$\to$} F,and r: F \hand{$\to$} $F_o$ is a left inverse(i.e. a retraction ),
and if \xout{ir} p = ir : F \hand{$\to$} F is in W,  so is r (and hence i too).

The condition c) \inserted{here} is the ingredient slightly stronger than what we used
to call "mild saturation property" of W,  meaning a) and b).

We'll call \uline{homotopy \xout{seg}interval in} \hand{$\mathcal{A}$} (with respect to the notion of
"weak equivalence" W,  which will generally be implicit and specified by
context)  a triple (I,$e_o$,$e_1$) ,  where I is an object of \hand{$\mathcal{A}$} , $e_o$ and $e_1$
\margin{* \unsure{provisional} \uline{separating} W-homotopy interval} two subobjects,  such that the following conditions a) to c) hold \hand{$(^{43})$}

a)  $e_o$ and $e_1$ are \xout{\illegible{1}} final objects of \hand{$\mathcal{A}$} ,
which implies that \hand{$\mathcal{A}$} has a final object, unique up to unique isomorphism,
which we'll debote by e$_{\hand{$\mathcal{A}$}}$ or simply e , so that the data  $e_o$,$e_1$ in I are
equivalent to giving two sections
\[ \hand{$\delta$}_o , \hand{$\delta$}_1 :\ e \hand{$\to$} I \]
of I over e.  Note that  $e_o$\hand{$\cap$}$e_1$ \hand{$=$} Ker(\hand{$\delta$}$_o$,\hand{$\delta$}$_1$).

b)  $e_o$\hand{$\cap$}$e_1$ = \hand{$\emptyset_{\mathcal{A}}$} (a strict initial element of \hand{$\mathcal{A}$} ) ,

c)  I \hand{$\to$} e is "universally in W" or,  as we'll say, is W-\uline{aspheric} or
simply aspheric ,

which just means here that \inserted{I $\to$ e is "squarable" and that} for any  F in Ob \hand{$\mathcal{A}$} ,  FxI \hand{$\to$} F is in W i.e.
is a "weak equivalence".

It is clear that if (I,$e_o$,$e_1$) is a homotopy \xout{segment}\inserted[typed]{interval} in \hand{$\mathcal{A}$}, then
for any F \hand{$\in$} Ob \hand{$\mathcal{A}$} ,  the"induced \xout{segment}\inserted[typed]{interval}"in \hand{$\mathcal{A}$}$_{/F}$ ,  namely (IxF,$e_o$xF,$e_1$xF)
is a homotopy interval in \hand{$\mathcal{A}$}$_{/F}$ .
